% Draft of Section 1, following planning/CMAME_PAPER_OUTLINE.md.
% Citations are verified keys from references/references_verified.bib.
% Figure 1 placeholder is marked below.
% TODO: replace hard-coded section numbers with \ref{} once all labels exist.
\section{Introduction}
\label{sec:intro}

\subsection{A Position-Only Shell from a Checkerboard Pattern}
\label{sec:intro-checkerboard}

Thin shells carry loads efficiently by combining membrane action with bending,
but their design and active deformation are complicated by the rotational
degrees of freedom carried by most shell elements. A position-only element would make the
geometry both the analysis state and the design variable; what is missing is a
discrete surface that supports a complete shell energy with positions alone.
This paper uses a checkerboard pattern for that purpose.

Consider a quadrilateral grid whose vertices are colored in a checkerboard
pattern. Discard the vertices of one color and keep the midpoints of the edges
of the remaining quads. Each quad is then replaced by the four midpoints of its
edges. These four points are coplanar and form a parallelogram, the Varignon
parallelogram, for any spatial position of the original quad. The resulting
surface is a collection of flat quadrilateral facets whose vertices are edge
midpoints of the original grid.

The construction is attractive for a shell element. The local tangent plane and
normal follow directly from nodal positions; no diagonal choice is required; and
the membrane metric and the bending curvature can share one geometric carrier.
The element developed here uses one position vector per node, three degrees of
freedom per node, and no rotational degrees of freedom.\footnote{Here
``Checkerboard'' refers to the midpoint pattern, not the checkerboard
instability of mixed finite elements.} The question is whether this minimal
state can carry a complete shell energy, and whether the resulting element is
useful when the shell geometry itself is the design variable.

% Figure 1 (to be inserted here): checkerboard pattern, edge midpoints, and the
% B/W patch. The figure should show the same pattern as a discrete surface and
% as a shell patch.

\subsection{Checkerboard Patterns in Geometry Processing}
\label{sec:intro-graphics}

Checkerboard and midpoint constructions have a long history in geometry processing. In discrete differential
geometry they appear as discrete isothermic surfaces and checkerboard nets
\cite{bobenkoPinkall1996,dellinger2024}. In graphics, position-based and
midedge shell models use edge midpoints to define bending
\cite{chenEtAl2018,chenEtAl2026}, while discrete shape operators and
subdivision shells build curvature from positions
\cite{grinspunEtAl2003,grinspunEtAl2006,cirakEtAl2000}. Smooth-basis
isogeometric shells take a different route: they carry exact geometry through
the basis \cite{hughesEtAl2005,kiendlEtAl2009}. The appeal of the midpoint
pattern is its locality, exact planarity, and the fact that one carrier serves
both membrane and bending. The pattern is attractive not only because it looks
right: the local frames and fundamental forms are available in closed form,
which makes it a convenient primitive for geometry-driven design.

Here we use the geometric primitive as the starting point for a mechanical formulation. The question
is mechanical: whether the midpoint surface can carry a reliable shell energy
under large deformation, and whether it can support design loops in which
geometry is the variable.

\subsection{Why Bring It into Mechanics}
\label{sec:intro-motivation}

The motivation goes beyond turning a geometric pattern into a shell. A position-only element would change how shell design is posed.
First, the state is the design variable: nodal positions are both the analysis
unknown and the geometric design variable; no rotational parameterization or
rotational constraint is needed, and the adjoint unknown has the same
three-degree-of-freedom structure as the state. Second, design loops become
cheaper: target-shape design, actuator layout, and shape optimization require
one forward and one adjoint solve per gradient instead of $2N$ forward solves.
Third, large rotations are handled by position updates: deployable and folding
structures do not need a rotation parameterization in the design loop. Fourth,
one geometric carrier serves membrane and bending: the same midpoint surface
provides both the metric and the curvature, without a separate director field.

These benefits point to applications such as morphing skins, active panels,
deployable booms and stents, soft actuators, and compliant mechanisms. They are,
however, potential benefits. A geometric pattern that looks natural can still
produce a discrete energy with spurious soft modes or a wrong thickness scaling.
Whether the benefits are realized depends on a mechanical question, which we
turn to next.

\subsection{From Discrete Geometry to Mechanics: The Gap}
\label{sec:intro-gap}

Geometry processing asks whether a discrete surface approximates a smooth one.
Mechanics asks a different question: whether the discrete energy is complete at
the lattice scale and in the thin-shell limit. Three issues matter. First, the
metric and the symmetric curvature are local averages; they can be consistent
to second order on smooth data and still be blind to lattice-scale modes.
Second, membrane and bending energies scale with thickness as $t$ and $t^3$
\cite{koiter1966,ciarlet2005}; a completion that does not respect this
separation can turn a bending mode into a membrane mode. Third, the energy
must remain objective under rigid rotation, especially when the design loop
explores large rotations \cite{simoFox1989}.

We call the modes that the midpoint metric and symmetric curvature cannot see
\emph{blind modes}. For the periodic midpoint surface, the analysis identifies
one membrane mode and two curvature modes. The migration from discrete geometry
to mechanics is therefore a formulation problem: identify these modes and
complete them without destroying the bending space. The central question is whether a minimal, targeted completion can do
that while preserving the simplicity of the position-only state. This is the
same tension that appears in membrane locking, where consistency of the strain
approximation is not enough and the discrete kernel matters
\cite{chapelleBathe2011,arnoldBrezzi1997,hiemstraEtAl2023}.

\subsection{This Paper's Approach}
\label{sec:intro-approach}

The paper follows four steps. We first define the midpoint surface, the flat
$B$ faces that carry the membrane metric, and the $W$-centered construction that
carries the second fundamental form. The same midpoint surface serves both;
boundary closure is local. We then analyze what the midpoint metric and
symmetric curvature miss. The membrane metric misses one alternating in-plane
mode; the symmetric curvature misses two transverse modes. We construct one
scalar completion for the membrane mode and one objective, reference-relative
curvature completion for the two curvature modes. Each completion acts only on
the mode it is designed for.

The nonlinear tangent inherits this structure. It is a sparse local matrix plus
one global rank-one term; the global term is solved with a bordered system and
is never assembled as a dense outer product. This makes the tangent exact for
the discrete energy and the design loop cheap. Finally, we use the element for
target-shape design: given a desired shape, solve for actuator forces or for a
low-dimensional set of initial geometry parameters. The examples are chosen to
expose both the capability and the boundary of the inverse problem.

\subsection{Contributions}
\label{sec:intro-contributions}

The contributions are fourfold. First, we construct a
three-degree-of-freedom-per-node shell element from the checkerboard midpoint
geometry; the state is nodal position only, and the membrane and bending
energies share one geometric carrier. Second, we identify the one membrane and
two curvature modes that the midpoint metric and symmetric curvature miss, and
give two targeted completions that act only on those modes. Third, we derive an
exact nonlinear tangent that is local and sparse apart from one global rank-one
term, and solve the global term with a bordered system. Fourth, we demonstrate
target-shape design: a nine-actuator shape-control problem reaches a
$0.0065\%$ target error, a two-parameter initial-curvature design through a
$90^\circ$ rotation recovers the target to $10^{-6}$, and the adjoint design
loop needs about $5.7$ times fewer linear solves than finite differences on the
same model.

\subsection{Scope and Preview}
\label{sec:intro-scope}

The element is aimed at design, active deformation, and large rotation. We
calibrate it on standard shell benchmarks and compare it with a midedge shell
\cite{chenEtAl2026} and MITC4 \cite{dvorkinBathe1984,batheDvorkin1985}. The
comparison establishes where the element sits, with the benchmarks used for
calibration. The operating range, including the mesh-thickness rule, is stated
in Section~4; the trade-offs are discussed in Section~7.

The remainder is organized as follows. Section~2 defines the element. Section~3
identifies the blind modes and constructs the completions. Section~4 presents
analytical properties and numerical implementation. Section~5 reports
verification and benchmarks. Section~6 presents the design applications.
Section~7 discusses the scope and concludes.
